\documentclass[12pt, a4paper]{article}

\usepackage[french]{babel}
\usepackage[utf-8]{inputenc}
\usepackage[left=2cm, right=2cm, top=2cm, bottom=2cm]{geometry}
\usepackage{graphicx}
\usepackage{listings}
\usepackage{xcolor}
\usepackage{hyperref}
\usepackage{fancyhdr}
\usepackage{tikz}
\usepackage{amsmath}
\usepackage{booktabs}
\usepackage{tcolorbox}

% Configuration listings
\lstset{
  basicstyle=\ttfamily\small,
  breaklines=true,
  frame=single,
  backgroundcolor=\color{gray!10},
  keywordstyle=\color{blue},
  commentstyle=\color{gray},
  stringstyle=\color{orange},
  showstringspaces=false,
  language=bash,
  morekeywords={docker, docker-compose}
}

% Configuration hyperref
\hypersetup{
  colorlinks=true,
  linkcolor=blue,
  urlcolor=blue
}

% En-tête et pied de page
\pagestyle{fancy}
\fancyhf{}
\fancyhead[L]{\textit{Séance 9 - Spark Structured Streaming}}
\fancyhead[R]{\thepage}
\fancyfoot[C]{\textit{Guide étudiant — earthquake\_streaming}}

\title{%
  \textbf{Spark Structured Streaming en Temps Réel} \\
  \vspace{0.5cm}
  \large Projet : Pipeline Tremblements de Terre (USGS) \\
  \vspace{0.3cm}
  \normalsize Séance 9 — Apache Spark
}
\author{}
\date{}

\begin{document}

\maketitle

\tableofcontents

\newpage

% ======================================================================
\section{Contexte et Objectifs}
% ======================================================================

Ce TP vous guide pour construire et exécuter un \textbf{pipeline temps réel} en trois couches :

\begin{enumerate}
  \item \textbf{Producer} — Interroge l'API USGS toutes les 30 secondes et produit des fichiers JSON.
  \item \textbf{Spark Structured Streaming} — Consomme les fichiers, enrichit les événements, écrit sur deux sinks en parallèle.
  \item \textbf{Orchestration} — Airflow déclenche et valide le pipeline.
\end{enumerate}

\subsection{Données}
Les données proviennent de l'API publique USGS (Earthquake Hazards Program) :
\begin{itemize}
  \item URL : \texttt{https://earthquake.usgs.gov/fdsnws/event/1/query}
  \item Format : GeoJSON
  \item Filtre : tremblements de magnitude $\geq 2.5$
  \item Fréquence : mise à jour toutes les 30 secondes
\end{itemize}

\subsection{Résultats attendus}
À la fin du TP, vous aurez :
\begin{itemize}
  \item Des fichiers \textbf{Parquet partitionnés} par date (stockage froid).
  \item Une table \textbf{PostgreSQL} avec tremblements enrichis (analytique).
  \item Un \textbf{DAG Airflow} orchestrant le tout.
\end{itemize}

% ======================================================================
\section{Architecture Générale}
% ======================================================================

\subsection{Flux de données}

\begin{center}
\begin{tikzpicture}[
  box/.style={rectangle, draw, minimum width=3cm, minimum height=0.8cm, text centered},
  arrow/.style={->, thick},
  label/.style={font=\small, text width=3cm, text centered}
]

% Boxes
\node[box] (usgs) at (0, 4) {USGS API};
\node[box] (producer) at (0, 2.5) {Producer \\ (Python)};
\node[box] (jsonl) at (0, 0.8) {/data/incoming \\ (JSONL files)};
\node[box] (spark) at (3, 2.5) {Spark Structured \\ Streaming};
\node[box] (parquet) at (5, 4) {Parquet \\ /data/clean/earthquakes};
\node[box] (postgres) at (5, 1) {PostgreSQL \\ earthquake\_enriched};

% Arrows
\draw[arrow] (usgs) -- (producer) node[midway, right] {poll 30s};
\draw[arrow] (producer) -- (jsonl) node[midway, right] {write atomic};
\draw[arrow] (jsonl) -- (spark) node[midway, above] {readStream};
\draw[arrow] (spark) -- (parquet) node[midway, above] {writeStream};
\draw[arrow] (spark) -- (postgres) node[midway, below] {foreachBatch};

\end{tikzpicture}
\end{center}

\subsection{Stack Docker}

Le projet fournit une composition Docker complète :

\begin{table}[h]
\centering
\small
\begin{tabular}{|l|l|l|}
\toprule
\textbf{Service} & \textbf{Image} & \textbf{Rôle} \\
\midrule
\texttt{data-init} & busybox & Crée les dossiers partagés \\
\texttt{postgres} & postgres:15 & Metadata + analytique \\
\texttt{spark-master} & bde2020/spark:3.3.0 & Master Standalone \\
\texttt{spark-worker} & bde2020/spark:3.3.0 & Worker (2 cores, 2 Go RAM) \\
\texttt{earthquake-producer} & python:3.11-slim & Producer USGS \\
\texttt{airflow-webserver} & custom & UI Airflow \\
\texttt{airflow-scheduler} & custom & Scheduler + driver Spark \\
\bottomrule
\end{tabular}
\end{table}

Tous les services partagent un volume Docker nommé \texttt{eq\_data} qui sert de « bus » entre le producer (producteur de fichiers) et Spark (consommateur).

% ======================================================================
\section{Composants Détaillés}
% ======================================================================

\subsection{Producer — \texttt{producer/earthquake\_producer.py}}

Le producer est une boucle infinie qui :

\begin{enumerate}
  \item Interroge l'API USGS avec une fenêtre glissante de 1 heure (par défaut).
  \item Déduplique les événements déjà vus (variable \texttt{seen\_ids}).
  \item Écrit atomiquement dans \texttt{/data/incoming/eq\_<YYYYMMDDTHHMMSS>.json}.
\end{enumerate}

\subsubsection{Écriture atomique}

\begin{tcolorbox}[colback=yellow!5, title=Astuce clé]
  \textbf{Spark file source indexe les fichiers dès qu'ils apparaissent.} \\
  Sans écriture atomique, le producer pourrait laisser un fichier partiellement écrit. Spark le lirait → crash.
\end{tcolorbox}

Le code écrit d'abord dans \texttt{eq\_<ts>.json.tmp}, puis utilise \texttt{os.replace()} pour renommer :

\begin{lstlisting}[language=python, caption=Écriture atomique]
with tmp.open("w", encoding="utf-8") as f:
    for row in rows:
        f.write(json.dumps(row) + "\n")
os.replace(tmp, final)  # atomic rename
\end{lstlisting}

\subsubsection{Format JSONL produit}

Chaque ligne du fichier est un événement sérialisé :

\begin{lstlisting}[language=json]
{"event_id": "us7000nabc", "magnitude": 4.2, "mag_type": "ml",
 "place": "10 km NW of Tokyo, Japan", "event_time": "2026-04-14T08:15:23+00:00",
 "latitude": 35.68, "longitude": 139.69, "depth_km": 12.3,
 "significance": 272, "tsunami": 0, "status": "reviewed",
 "type": "earthquake", "title": "M 4.2 - 10 km NW of Tokyo",
 "ingested_at": "2026-04-14T08:16:00+00:00"}
\end{lstlisting}

\subsection{Job Spark — \texttt{spark\_jobs/earthquake\_streaming.py}}

\subsubsection{Lecture (readStream)}

\begin{lstlisting}[language=python, caption=Lecture streaming]
spark.readStream \
    .schema(EVENT_SCHEMA) \
    .option("maxFilesPerTrigger", 5) \
    .json(INCOMING_DIR)
\end{lstlisting}

\begin{itemize}
  \item \texttt{schema} explicite : plus rapide que l'inférence, robuste aux fichiers mal formés.
  \item \texttt{maxFilesPerTrigger=5} : chaque trigger traite au maximum 5 fichiers.
  \item Trigger par défaut : toutes les 10 secondes.
\end{itemize}

\subsubsection{Enrichissement}

Quatre colonnes sont ajoutées à chaque événement :

\begin{table}[h]
\centering
\small
\begin{tabular}{|l|l|}
\toprule
\textbf{Colonne} & \textbf{Règle} \\
\midrule
\texttt{depth\_category} & $<70$ km = Superficiel, $<300$ = Intermédiaire, sinon Profond \\
\texttt{severity} & tsunami=1 → TSUNAMI, mag $\geq 8$ → Dévastateur, etc. \\
\texttt{is\_significant} & 1 si \texttt{significance} $> 500$, sinon 0 \\
\texttt{event\_date} & Date formatée pour partitionnement \\
\bottomrule
\end{tabular}
\end{table}

\subsubsection{Deux sinks en parallèle}

Le job écrit sur \textbf{deux cibles différentes} via un seul \texttt{foreachBatch} :

\begin{enumerate}
  \item \textbf{Parquet} : partitionné par \texttt{event\_date}, stockage froid.
  \item \textbf{PostgreSQL} : table \texttt{earthquake\_enriched}, analytique.
\end{enumerate}

\begin{tcolorbox}[colback=yellow!5, title=Important]
  \textbf{Ne pas utiliser deux \texttt{writeStream} distincts !} \\
  Ils partageraient le sourceLog et causeraient : \texttt{IllegalStateException: Concurrent update to the log}.
\end{tcolorbox}

\subsection{DAG Airflow — \texttt{dags/earthquake\_streaming\_dag.py}}

La DAG a quatre tâches :

\begin{lstlisting}
health_check --> init_postgres --> +-> submit_spark_streaming (long-running)
                                   |
                                   +-> validate_run (check apres 3 min)
\end{lstlisting}

\subsubsection{health\_check (BashOperator)}

Vérifie que les trois dossiers partagés existent et sont accessibles :

\begin{lstlisting}[language=bash]
mkdir -p /data/incoming /data/clean/earthquakes /data/checkpoints/earthquake
ls -ld /data/incoming /data/clean/earthquakes /data/checkpoints/earthquake
\end{lstlisting}

\subsubsection{init\_postgres (PostgresOperator)}

Crée la table de destination (idempotente) :

\begin{lstlisting}[language=sql]
CREATE TABLE IF NOT EXISTS earthquake_enriched (
    event_id       TEXT PRIMARY KEY,
    magnitude      DOUBLE PRECISION,
    ...
    processed_at   TIMESTAMPTZ
);
\end{lstlisting}

\subsubsection{submit\_spark\_streaming (SparkSubmitOperator)}

Lance le job Spark en mode client. Le driver s'exécute dans Airflow, les executors sur le worker Spark.

\textbf{Important :} cette tâche est long-running et reste en état \texttt{running} tant que le job Spark tourne.

\subsubsection{validate\_run (PythonOperator)}

En parallèle de \texttt{submit\_spark\_streaming}, cette tâche :
\begin{enumerate}
  \item Attend 3 minutes (le temps d'un producer poll + un trigger Spark).
  \item Compte les lignes en PostgreSQL et les fichiers Parquet.
  \item Log un rapport de validation.
  \item Échoue si zéro lignes ET zéro fichiers (pipeline non productif).
\end{enumerate}

% ======================================================================
\section{Démarrage Étape par Étape}
% ======================================================================

\subsection{Prérequis}

\begin{itemize}
  \item Docker et docker-compose installés.
  \item Au minimum 8 Go de RAM disponibles (4 pour Spark + 2 pour Airflow + 2 pour Postgres).
  \item Connexion Internet (le producer interroge l'API USGS).
\end{itemize}

\subsection{Étape 1 : Construire les images Docker}

Allez dans le dossier \texttt{infra/} :

\begin{lstlisting}
cd infra/
docker compose build
\end{lstlisting}

Cela construit l'image personnalisée \texttt{Dockerfile.airflow} qui inclut PySpark et le driver JDBC PostgreSQL.

Durée attendue : 2--3 minutes.

\subsection{Étape 2 : Démarrer la stack}

\begin{lstlisting}
docker compose up -d
\end{lstlisting}

Les conteneurs démarrent dans cet ordre (via \texttt{depends\_on}) :
\begin{enumerate}
  \item \texttt{data-init} crée les dossiers, puis termine.
  \item \texttt{postgres} démarre.
  \item \texttt{spark-master} et \texttt{spark-worker} démarrent en parallèle.
  \item \texttt{earthquake-producer} démarre et entre sa boucle infinie.
  \item \texttt{airflow-init} migre la base, puis termine.
  \item \texttt{airflow-webserver} et \texttt{airflow-scheduler} démarrent.
\end{enumerate}

Vérifiez l'état :

\begin{lstlisting}
docker compose ps
\end{lstlisting}

Attendez que tous les services soient \texttt{running} et \texttt{healthy}.

\subsection{Étape 3 : Vérifier que le producer produit}

\begin{lstlisting}
docker compose logs -f earthquake-producer | head -20
\end{lstlisting}

Vous devriez voir des logs comme :

\begin{lstlisting}
2026-04-14 08:15:34,123 [INFO] Fetched 147 features (08:13:00 -> 08:15:34)
- 42 new, 105 already seen
2026-04-14 08:15:34,456 [INFO] Wrote 42 new rows ->
/data/incoming/eq_20260414T081534.json
\end{lstlisting}

Vérifiez les fichiers :

\begin{lstlisting}
docker exec -it eq-spark-worker ls -lh /data/incoming/ | head -5
\end{lstlisting}

Vous devriez voir plusieurs fichiers \texttt{eq\_*.json}.

\subsection{Étape 4 : Accéder à Airflow}

Ouvrez votre navigateur à \texttt{http://localhost:8085}.

Login : \texttt{admin} / \texttt{admin}

Vous devriez voir le DAG \texttt{EARTHQUAKE-STREAMING} dans la liste.

\subsection{Étape 5 : Déclencher le DAG}

Cliquez sur le DAG, puis sur \texttt{Trigger DAG}.

Le scheduler crée un \texttt{dag\_run} et exécute les tâches.

Attendez que :
\begin{itemize}
  \item \texttt{health\_check} soit \texttt{success}.
  \item \texttt{init\_postgres} soit \texttt{success}.
  \item \texttt{submit\_spark\_streaming} passe à \texttt{running}.
  \item \texttt{validate\_run} commence à s'exécuter (elle attend 3 min).
\end{itemize}

\subsection{Étape 6 : Attendre la validation}

Attendez 3--4 minutes. \texttt{validate\_run} va :
\begin{enumerate}
  \item Interroger PostgreSQL.
  \item Compter les fichiers Parquet.
  \item Afficher un rapport.
\end{enumerate}

Consultez les logs de la tâche dans Airflow pour confirmer le rapport.

% ======================================================================
\section{Vérifier les Résultats}
% ======================================================================

\subsection{1. Vérifier le producer}

\begin{lstlisting}[caption=Consulter les logs du producer]
docker compose logs --tail 50 earthquake-producer
\end{lstlisting}

Cherchez des lignes comme :

\begin{lstlisting}
[INFO] Wrote 42 new rows -> /data/incoming/eq_20260414T081534.json
\end{lstlisting}

\begin{lstlisting}[caption=Lister les fichiers produits]
docker exec -it eq-spark-worker ls -lh /data/incoming/
\end{lstlisting}

\subsection{2. Vérifier Parquet}

\begin{lstlisting}[caption=Lister la structure de Parquet]
docker exec -it eq-spark-worker ls -R /data/clean/earthquakes/
\end{lstlisting}

Vous devriez voir une structure comme :

\begin{lstlisting}
/data/clean/earthquakes/
event_date=2026-04-14/
  part-00000-xyz.parquet
event_date=2026-04-13/
  part-00000-abc.parquet
\end{lstlisting}

\begin{lstlisting}[caption=Lire les données Parquet avec Spark]
docker exec -it eq-spark-worker spark-shell --master local[*] \
    -e "spark.read.parquet(\"/data/clean/earthquakes\").show(10, false)"
\end{lstlisting}

\subsection{3. Vérifier PostgreSQL}

\begin{lstlisting}[caption=Se connecter à PostgreSQL]
docker exec -it eq-postgres psql -U airflow -d earthquakes
\end{lstlisting}

Puis exécutez :

\begin{lstlisting}[language=sql]
-- Top 10 événements récents
SELECT event_id, magnitude, place, severity, processed_at
FROM earthquake_enriched
ORDER BY event_ts DESC LIMIT 10;

-- Distribution par sévérité
SELECT severity, COUNT(*), ROUND(AVG(magnitude)::numeric, 2) AS avg_mag
FROM earthquake_enriched
GROUP BY severity
ORDER BY severity DESC;
\end{lstlisting}

\subsection{4. Spark UI}

Accédez à \texttt{http://localhost:8080} pour voir la Spark Master UI.

Cliquez sur l'application \texttt{earthquake-structured-streaming}, puis allez dans l'onglet \textbf{Structured Streaming} pour voir :
\begin{itemize}
  \item Input rate (lignes/seconde lues).
  \item Processing rate (lignes/seconde traitées).
  \item Latence des micro-batches.
\end{itemize}

% ======================================================================
\section{Variables d'Environnement (Tuning)}
% ======================================================================

Toutes ces variables peuvent être modifiées dans \texttt{infra/.env.airflow} avant de lancer \texttt{docker compose up} :

\begin{table}[h]
\centering
\small
\begin{tabular}{|l|l|l|}
\toprule
\textbf{Variable} & \textbf{Défaut} & \textbf{Description} \\
\midrule
\texttt{POLL\_INTERVAL} & 30 & Secondes entre deux polls USGS \\
\texttt{MIN\_MAGNITUDE} & 2.5 & Magnitude minimale (filtre USGS) \\
\texttt{INITIAL\_LOOKBACK\_DAYS} & 7 & Lookback au premier poll \\
\texttt{POLL\_WINDOW\_SECONDS} & 3600 & Fenêtre glissante en secondes \\
\midrule
\texttt{INCOMING\_DIR} & /data/incoming & Répertoire producer \\
\texttt{CLEAN\_DIR} & /data/clean/earthquakes & Répertoire Parquet \\
\texttt{CHECKPOINT\_ROOT} & /data/checkpoints/earthquake & Répertoire checkpoints \\
\midrule
\texttt{PG\_URL} & jdbc:postgresql://postgres:5432/earthquakes & URL JDBC \\
\texttt{PG\_USER} & airflow & User Postgres \\
\texttt{PG\_PASSWORD} & airflow & Password Postgres \\
\bottomrule
\end{tabular}
\end{table}

\subsection{Exemple : Augmenter la fréquence de traitement}

\begin{lstlisting}[caption=Éditer .env.airflow, language=bash]
# Réduire le poll interval a 15 secondes
POLL_INTERVAL=15

# Augmenter le trigger Spark: editer spark_jobs/earthquake_streaming.py
# trigger(processingTime="5 seconds")  # au lieu de 10
\end{lstlisting}

Puis redémarrer :

\begin{lstlisting}
docker compose down
docker compose up -d
\end{lstlisting}

% ======================================================================
\section{Points Techniques Notables}
% ======================================================================

\subsection{Concurrence sur le FileStreamSource}

\begin{tcolorbox}[colback=red!5, title=Piège courant]
  \textbf{Ne jamais utiliser deux \texttt{writeStream} sur la même source !}

  Ils partagent un sourceLog interne. Chacun essaie de mettre à jour les offsets → race condition → crash.

  \textit{Solution :} utiliser un seul \texttt{writeStream} avec \texttt{foreachBatch} pour écrire sur plusieurs cibles.
\end{tcolorbox}

\subsection{Permissions cross-UID sur volume Docker}

\begin{tcolorbox}[colback=orange!5, title=Détail d'implémentation]
  Le producer tourne en uid 0 (root) dans \texttt{spark-worker}. Airflow tourne en uid 50000 dans \texttt{airflow-scheduler}. Spark Structured Streaming crée des répertoires temporaires lors du commit → problème de droits.

  \textit{Corrections :}
  \begin{itemize}
    \item OS-level : \texttt{umask 000} avant lancer spark-master/worker.
    \item Hadoop-level : \texttt{spark.hadoop.fs.permissions.umask-mode=000} en config Spark.
  \end{itemize}
\end{tcolorbox}

\subsection{Checkpoints séparés pour chaque sink}

\begin{tcolorbox}[colback=blue!5, title=Bonne pratique]
  Chaque query (même si elles partagent une source) doit avoir son propre checkpoint.

  Ici : \texttt{/data/checkpoints/earthquake/combined} est utilisé pour le foreachBatch.

  Si vous aviez deux writeStream historiquement, vous auriez \texttt{/data/checkpoints/earthquake/parquet} et \texttt{/data/checkpoints/earthquake/postgres}. \textbf{Ne jamais les mutualiser.}
\end{tcolorbox}

\subsection{Déduplication côté producer}

Le producer maintient \texttt{seen\_ids} en mémoire. Entre deux redémarrages du producer, cette set est perdue. Donc, si le producer est redémarré et repolled avec la même fenêtre, il peut re-produire des événements.

\textit{Mitigation en Postgres :} la table a une PK sur \texttt{event\_id} et Spark utilise \texttt{ON CONFLICT DO NOTHING} → les doublons sont ignorés.

% ======================================================================
\section{Dépannage}
% ======================================================================

\subsection{« Producer n'écrit rien »}

\begin{lstlisting}
docker compose logs earthquake-producer | tail -30
\end{lstlisting}

Cherchez \texttt{RequestException} ou \texttt{timeout}. Vérifiez votre connexion Internet. L'API USGS peut être lente.

\textit{Solution :} augmenter \texttt{timeout=15} dans \texttt{producer/earthquake\_producer.py} si le réseau est lent.

\subsection{« Spark n'écrit rien »}

\begin{lstlisting}
docker compose logs airflow-scheduler | tail -50
\end{lstlisting}

Cherchez \texttt{ERROR} ou \texttt{Exception}. Vérifiez :
\begin{itemize}
  \item Les dossiers \texttt{/data/incoming}, \texttt{/data/clean/earthquakes}, \texttt{/data/checkpoints/earthquake} existent-ils ?
  \item \texttt{health\_check} est-elle passée ?
  \item Spark peut-il lire les fichiers JSONL : vérifiez le schéma.
\end{itemize}

\subsection{« Postgres est vide »}

\begin{lstlisting}
docker exec -it eq-postgres psql -U airflow -d earthquakes \
    -c "SELECT COUNT(*) FROM earthquake_enriched;"
\end{lstlisting}

Si zéro lignes :
\begin{itemize}
  \item Vérifiez \texttt{validate\_run} : a-t-elle échoué ou attendu ?
  \item Vérifiez les logs du job Spark pour les erreurs JDBC.
  \item Testez la connexion JDBC manuellement.
\end{itemize}

\subsection{« Airflow : tâche stuck en running »}

\texttt{submit\_spark\_streaming} est censée rester en \texttt{running} (c'est une query long-running Spark).

Pour \textbf{arrêter proprement} le DAG :
\begin{enumerate}
  \item Dans l'UI Airflow, cliquez sur \texttt{submit\_spark\_streaming} → \texttt{Terminate}.
  \item Ou : \texttt{docker compose down} (arrête tous les conteneurs).
\end{enumerate}

\subsection{« Volume Docker ne se monte pas »}

\begin{lstlisting}
docker compose down -v
docker compose up -d
\end{lstlisting}

Cela réinitialise les volumes. \texttt{data-init} recrée les dossiers.

% ======================================================================
\section{Arrêt et Reset}
% ======================================================================

\subsection{Arrêt propre (garde les données)}

\begin{lstlisting}
cd infra/
docker compose down
\end{lstlisting}

Les checkpoints, fichiers Parquet, et données Postgres sont \textbf{préservés}. Vous pouvez relancer \texttt{docker compose up -d} et Spark reprendra \emph{exactly-once} où il s'était arrêté.

\subsection{Réinitialisation complète}

\begin{lstlisting}
cd infra/
docker compose down -v
\end{lstlisting}

Cela supprime \textbf{tous les volumes} (\texttt{eq\_data}, \texttt{eq\_postgres\_data}, \texttt{eq\_airflow\_logs}). À partir de maintenant, tout repart de zéro.

\subsection{Rebuild et redémarrage}

Si vous modifiez un fichier Python (\texttt{producer}, job Spark, DAG), redémarrez comme suit :

\begin{lstlisting}
docker compose down
docker compose build
docker compose up -d
\end{lstlisting}

% ======================================================================
\section{Extensions Possibles}
% ======================================================================

Une fois que vous maîtrisez le pipeline, vous pouvez expérimenter :

\begin{enumerate}
  \item \textbf{Kafka au lieu de file source} : remplacer \texttt{readStream.json(INCOMING\_DIR)} par \texttt{readStream.format("kafka")}.
  \item \textbf{Alertes en temps réel} : ajouter une query pour déterminer les événements « significatifs » et les pousser vers une queue externe.
  \item \textbf{Prédictions ML} : charger un modèle scikit-learn ou MLlib, scorer chaque événement.
  \item \textbf{Backfill} : relancer le job avec une lookback de plusieurs mois pour re-traiter l'historique.
\end{enumerate}

% ======================================================================
\section{Conclusion}
% ======================================================================

Ce TP illustre les concepts clés de \textbf{Spark Structured Streaming} :

\begin{itemize}
  \item \textbf{readStream} sur une source.
  \item \textbf{Enrichissement} et calculs stateful.
  \item \textbf{Deux sinks} simultanés (Parquet + Postgres).
  \item \textbf{Checkpoints} pour la tolérance aux pannes.
  \item \textbf{Orchestration} via Airflow.
\end{itemize}

Le pattern « file source » est simplifié (pas de Kafka), idéal pour pédagogie. En production, vous remplaceriez par Kafka pour la durabilité et la scalabilité.

\vspace{1cm}

\textit{Bon apprentissage !}

\end{document}
